\section{Deep-Supervised VICReg Architecture}
\label{sec:architecture}

The current winning model is intentionally close to the stabilized
\CodeWM{} baseline.  It changes where supervision is applied, how far
the encoder is allowed to recurse, and how the final readout is
regularized; it does not rely on a larger model or a new dataset.

\subsection{Baseline Components}

The base model has three components:
\begin{enumerate}
    \item an online looped transformer encoder for the pre-edit code
          state;
    \item a lightweight action encoder for the edit vector;
    \item a looped predictor that maps online state plus action into the
          target encoder's post-edit space.
\end{enumerate}
The target encoder remains near-static with EMA decay $0.99999$, matching
the stabilization recipe from the first paper.

\subsection{\code{loops3-auxpred-vicreg}}

The current Phase~9 recipe is:
\begin{center}
\begin{tabular}{ll}
\toprule
Setting & Value \\
\midrule
Encoder loops & $3$ \\
Model dimension & $128$ \\
Attention heads & $4$ \\
Pooling & attention/CLS-style pooled code state \\
Predictor & 6 looped blocks, depth 2 \\
Target update & EMA decay $0.99999$ \\
Auxiliary loops & $1,2$ \\
Auxiliary type & predictive \\
Auxiliary weight & $0.3$ \\
Readout regularizer & VICReg \\
Signal regularizer weight & $0.1$ \\
\bottomrule
\end{tabular}
\end{center}

In environment-variable form:
\begin{center}
\begin{tabular}{ll}
\code{WM\_ENCODER\_LOOPS} & \code{3} \\
\code{WM\_AUX\_LOOPS} & \code{1,2} \\
\code{WM\_LAMBDA\_AUX} & \code{0.3} \\
\code{WM\_AUX\_TYPE} & \code{pred} \\
\code{WM\_REG\_MODE} & \code{vicreg} \\
\code{WM\_SIGREG\_WEIGHT} & \code{0.1} \\
\code{WM\_EMA\_DECAY} & \code{0.99999}
\end{tabular}
\end{center}

\subsection{Why Predictive Aux Beats Contrastive Aux}

The current evidence favors predictive auxiliary losses over contrastive
auxiliary losses for the loop-intermediate objective.  Contrastive aux
losses impose representation-level geometry at intermediate loops, while
predictive aux losses ask each loop to remain useful for the actual
online-to-target mapping.  For a looped encoder with shared weights, this
matters: the same block must be useful after one application, two
applications, and the final application.  Predictive deep supervision
directly trains that property.

\subsection{Why VICReg Is Needed at the Readout}

The high-lift \code{loops3-auxpred} checkpoint demonstrated that deep
supervision repaired the intermediates, not the whole representation
pipeline.  Per-loop probes still found strong edit information at loop 3,
but the final attention-pool readout collapsed to an online effective
rank of about $2/128$.  In that state the predictor can learn a real
online-to-target mapping while still failing to distinguish examples:
the predicted vector has high cosine with the correct target, but also
high cosine with many wrong targets.

VICReg targets this failure directly.  The invariance term keeps the
prediction aligned with the target, while the variance hinge and
covariance penalty prevent the pooled code state from using only a tiny
number of dimensions.  This is why the current recipe reports
discriminability metrics next to lift: a viable readout must improve both
prediction alignment and cross-example separation.

\subsection{Interpreting Negative \copycos}

In medium-tier \code{loops3-auxpred} runs, \copycos{} becomes
approximately zero or slightly negative.  This can mean the predictor is
doing real work by translating from online representation space into
target representation space.  It is not, by itself, enough to claim a
healthy representation.  The Phase~9 caveat is that a collapsed readout
can also make the copy baseline weak and thereby inflate lift.  Negative
\copycos{} is therefore useful only when paired with effective-rank and
KNN checks.
